\documentclass[11pt,a4paper,oneside]{report}

\usepackage[utf8]{inputenc}
\usepackage[T1]{fontenc}
\usepackage[brazilian]{babel}
\usepackage{lmodern}
\usepackage[a4paper,margin=2.5cm]{geometry}
\usepackage{graphicx}
\usepackage{float}
\usepackage{pdflscape}
\usepackage{xcolor}
\usepackage{microtype}
\usepackage{booktabs}
\usepackage{titlesec}
\usepackage{enumitem}
\usepackage{placeins}
\usepackage{afterpage}
\usepackage{hyperref}

\hypersetup{
  colorlinks=true,
  linkcolor=black!70,
  urlcolor=blue!60!black,
  citecolor=blue!60!black,
  pdftitle={Void Descent --- Documentacao Tecnica},
  pdfauthor={Equipe Void Descent},
  pdfsubject={Documentacao Tecnica},
  pdfkeywords={roguelike, webgl, electron, procedural, javascript}
}

% Cabecalho de capitulo mais compacto (sem espaco enorme acima)
\titleformat{\chapter}[hang]
  {\normalfont\huge\bfseries}
  {\thechapter.}{1ex}{}
\titlespacing*{\chapter}{0pt}{0pt}{20pt}

% Helper para diagramas em retrato com altura limitada
\newcommand{\diagfig}[3]{%
  \begin{figure}[ht!]
    \centering
    \includegraphics[width=\linewidth,height=0.78\textheight,keepaspectratio]{#1}
    \caption{#2}
    \label{fig:#3}
  \end{figure}%
  \FloatBarrier
}

% Helper para diagramas largos em paisagem (usa afterpage para evitar orphans)
\newcommand{\diagfiglandscape}[3]{%
  \afterpage{%
    \clearpage
    \begin{landscape}
    \begin{figure}[H]
      \centering
      \includegraphics[width=\linewidth,height=0.85\textheight,keepaspectratio]{#1}
      \caption{#2}
      \label{fig:#3}
    \end{figure}
    \end{landscape}
  }%
}

\graphicspath{{./}}

\begin{document}

% ====== Capa ======
\begin{titlepage}
  \centering
  \vspace*{3cm}

  {\Huge\bfseries VOID DESCENT}\\[0.6cm]
  {\Large Documentação Técnica}\\[0.5cm]
  \rule{0.6\textwidth}{0.4pt}\\[1cm]
  {\large Roguelike top-down procedural com WebGL}

  \vspace{2.5cm}

  {\large\textbf{Equipe}}\\[0.5cm]
  \begin{tabular}{c}
    Luis Gustavo Olímpio \\
    Lauan Matheus \\
    José Melquíades \\
    João Leonardi
  \end{tabular}

  \vfill

  iCEV --- Instituto de Ensino Superior\\
  Introdução ao Desenvolvimento de Jogos\\
  \vspace{0.3cm}
  Maio de 2026
\end{titlepage}

% ====== Sumario e Lista ======
\tableofcontents
\listoffigures

\clearpage

% \input nao forca \clearpage como \include faz, evitando paginas em branco
\chapter{Visão Geral}

\section{O jogo}

\textbf{Void Descent} é um \emph{roguelike} top-down em tempo real construído em
JavaScript puro com renderização WebGL. Cada partida regenera todo o conteúdo do jogo
(mapas, atlas de tiles, áudio, comportamento de inimigos) por algoritmos em tempo de
execução, sem nenhum \emph{asset} pré-fabricado.

A premissa é descer três andares procedurais, \textit{Antecâmara}, \textit{Subsistemas}
e \textit{Núcleo}, eliminando inimigos sala por sala e coletando armas e
\textit{upgrades}, até derrotar o boss final. A morte é permanente; cada \emph{run} é
uma nova execução do gerador.

\section{Equipe}

\begin{itemize}
  \item Luis Gustavo Olímpio
  \item Lauan Matheus
  \item José Melquíades
  \item João Leonardi
\end{itemize}

\section{Stack tecnológico}

\begin{itemize}
  \item \textbf{JavaScript ES6+}, sem \textit{frameworks}, ESLint ou bundlers.
    Carregamento direto de \texttt{<script>} em ordem.
  \item \textbf{WebGL 1.0} com \textit{shaders} de vértice e fragmento personalizados, e
    iluminação por vértice de até 32 luzes pontuais simultâneas.
  \item \textbf{Sonant-X}, \textit{port} JavaScript do sintetizador Sonant
    (Marcus Geelnard, Nicolas Vanhoren), gera todo o áudio em tempo de \emph{boot}.
  \item \textbf{Web Audio API}: \texttt{AudioContext} e \texttt{GainNode} mestre para
    \textit{routing} unificado e controle de \emph{mute}.
  \item \textbf{Electron 32} para empacotamento como \texttt{.exe} \textit{portable}
    Windows x64, via \texttt{electron-builder}.
\end{itemize}

\section{Métricas do projeto}

\begin{tabular}{lr}
  \toprule
  \textbf{Métrica} & \textbf{Valor} \\
  \midrule
  Linhas de código (JS+HTML) & $\sim$2.100 \\
  Tamanho do código fonte & 77 KB \\
  Quantidade de arquivos JS & 8 \\
  Tamanho do \texttt{.exe} \textit{portable} & 73 MB \\
  \quad Chromium embutido & $\sim$60 MB \\
  \quad Node runtime & $\sim$10 MB \\
  \quad Bundle do jogo (asar) & $\sim$80 KB \\
  Tempo médio de \emph{run} & 8 a 12 minutos \\
  \bottomrule
\end{tabular}

O contraste entre o código fonte (77 KB) e o \texttt{.exe} final (73 MB) decorre da
escolha de empacotamento via Electron, conforme especificação do GDD para garantir
experiência \emph{nativa} de \textit{desktop}. A geração procedural reside na camada do
\emph{bundle}, não no \emph{wrapper} de distribuição.

\section{Estrutura desta documentação}

A documentação está dividida em nove capítulos, cada um focando uma dimensão ortogonal
do projeto:

\begin{enumerate}
  \item \textbf{Visão Geral}, este capítulo.
  \item \textbf{Arquitetura de Componentes}, com a divisão modular e as dependências
    (\textit{component diagram}).
  \item \textbf{Sistema de Entidades}, com a hierarquia de classes do mundo do jogo
    (\textit{class diagram}).
  \item \textbf{Máquina de Estados}, com as transições do estado global
    (\textit{state diagram}).
  \item \textbf{Game Loop} \textit{micro} (sala-a-sala) e \textit{macro} (run completa),
    representado por \textit{activity diagrams}.
  \item \textbf{Sistemas Procedurais}: geração de \emph{dungeon}, síntese de áudio e
    \textit{tinting} de atlas (\textit{sequence diagrams}).
  \item \textbf{Build e Deploy}, com o \emph{pipeline} de empacotamento e distribuição
    (\textit{deployment diagram}).
  \item \textbf{Backlog Técnico, Decisões e Aprendizados}, com o caminho real do
    projeto: o que escolhemos, o que descartamos e o que aprendemos.
  \item \textbf{Considerações Finais}, com o balanço do que foi entregue, o que
    faríamos diferente e o que vem depois.
\end{enumerate}

\chapter{Arquitetura de Componentes}

\section{Decomposição modular}

O código está dividido em oito módulos JavaScript, cada um com responsabilidade única e
bem delimitada. O \textbf{Component Diagram} é o tipo UML adequado para representar a
divisão de \emph{software} em unidades \emph{deployáveis} e suas dependências de uso.

\diagfiglandscape{diagrams/components.png}{Diagrama de Componentes, com os módulos do Void Descent e suas dependências}{components}

\section{Responsabilidades por módulo}

\begin{description}

  \item[\texttt{src/vd/random.js}] Gerador de números pseudoaleatórios determinístico,
    baseado em \emph{Linear Congruential Generator} (LCG) de 32 bits. É \emph{seedado}
    explicitamente a cada andar (\texttt{random\_seed(0xBADC0DE1 + floor\_id)}), o que
    garante reprodutibilidade do mapa para uma dada combinação de \emph{seed} e andar.

  \item[\texttt{src/vd/sonantx.js}] \emph{Port} JavaScript do sintetizador Sonant.
    Implementa osciladores (sin, quadrada, dente-de-serra, triangular), envelopes ADSR,
    filtros (lopass, hipass, bandpass, notch), LFO e \emph{delay}, todos rodando sobre
    \texttt{Float32Array}.

  \item[\texttt{src/vd/audio.js}] Define os \emph{patches} de instrumentos (oito SFX e a
    faixa ambiente). Inicializa o \texttt{AudioContext}, cria o \texttt{GainNode} mestre
    (\texttt{audio\_master\_gain}) que centraliza o controle de \emph{mute}, e coordena
    a geração assíncrona dos \texttt{AudioBuffer}s no \emph{boot}.

  \item[\texttt{src/vd/renderer.js}] \textit{Pipeline} WebGL. Compila os \emph{shaders}
    de vértice e fragmento, gerencia o atlas de tiles (\texttt{renderer\_bind\_image}),
    expõe primitivas de geometria (\texttt{push\_quad}, \texttt{push\_floor},
    \texttt{push\_block}, \texttt{push\_sprite}) e gerencia até 32 luzes pontuais por
    \emph{frame}.

  \item[\texttt{src/vd/dungeon.js}] Algoritmo de \emph{room placement} com validação
    posterior. Posiciona retângulos não sobrepostos, conecta-os por corredores em forma
    de L, marca tipos especiais (\texttt{start}, \texttt{chest}, \texttt{boss}) e
    converte o resultado em geometria via chamadas a
    \texttt{push\_block}/\texttt{push\_floor}.

  \item[\texttt{src/vd/entities.js}] Hierarquia de classes do mundo simulado:
    classe-base \texttt{entity\_t}, \emph{player}, quatro tipos de inimigos,
    \emph{boss}, projéteis, partículas, explosões e \emph{pickups}. Detalhada no
    Capítulo~\ref{ch:entities}.

  \item[\texttt{src/vd/terminal.js}] \emph{Overlay} HUD textual no estilo terminal de
    boot. Renderiza no elemento DOM \texttt{<code id="a">} (fora do \emph{canvas}) e
    exibe avisos temporários como ``ANDAR 1 :: ANTECAMARA'' ou ``PORTAS TRANCADAS''.

  \item[\texttt{src/vd/game.js}] Módulo orquestrador. Mantém o \texttt{game\_state}
    global, implementa o \texttt{game\_tick} (despachado por estado), e gerencia a
    \emph{shop} de \emph{upgrades}, o \emph{lock-and-clear} de salas, o \emph{spatial
    grid} de colisão e a HUD principal.

\end{description}

\section{Dependências e fluxo de controle}

A arquitetura segue um padrão estritamente \textbf{em camadas}, sem ciclos:

\begin{enumerate}
  \item \textbf{Camada de utilitários puros}, sem dependências:
    \texttt{random.js}, \texttt{sonantx.js}.

  \item \textbf{Camada de infraestrutura}: \texttt{renderer.js}, \texttt{audio.js} (que
    depende de \texttt{sonantx.js}), \texttt{terminal.js}.

  \item \textbf{Camada de geração}: \texttt{dungeon.js}, que depende de
    \texttt{random.js} e da API de geometria de \texttt{renderer.js}.

  \item \textbf{Camada de simulação}: \texttt{entities.js}, que depende das três
    anteriores via globais.

  \item \textbf{Orquestração}: \texttt{game.js}, que depende de tudo.
\end{enumerate}

A ausência de um sistema de módulos (CommonJS, ESM) é intencional. Todos os símbolos
vivem no escopo global definido pela ordem de inclusão dos \texttt{<script>} em
\texttt{index.html}. O padrão minimalista é coerente com a proposta procedural do
projeto e elimina o custo de \emph{boot} de um \emph{bundler}.

\section{Boot sequence}

A inicialização do programa, definida em \texttt{index.html}, segue uma ordem rígida:

\begin{enumerate}
  \item \texttt{audio\_init()} cria o \texttt{AudioContext} e dispara a síntese
    assíncrona da música e dos SFX.
  \item \texttt{renderer\_init()} cria o contexto WebGL, compila os \emph{shaders} e
    aloca o \emph{vertex buffer}.
  \item \texttt{bind\_input()} registra os \emph{handlers} de teclado e mouse no
    \texttt{document}.
  \item \texttt{show\_menu()} define \texttt{game\_state = 'menu'} e instala os
    \emph{handlers} específicos da tela de título.
  \item \texttt{requestAnimationFrame(game\_tick)} inicia o \emph{loop} principal.
\end{enumerate}

\chapter{Sistema de Entidades}
\label{ch:entities}

\section{Modelo orientado a objetos}

Toda forma de vida (e morte) no mundo do Void Descent é uma \texttt{entity\_t}, uma
classe-base ECMAScript 2015 com herança simples. O \textbf{Class Diagram} comunica
simultaneamente herança, atributos, métodos e relacionamentos.

\diagfiglandscape{diagrams/entities-class.png}{Diagrama de Classes, com a hierarquia de \texttt{entity\_t}}{entities-class}

\section{A classe-base \texttt{entity\_t}}

Todas as entidades compartilham:

\begin{itemize}
  \item \textbf{Cinemática}: posição (\texttt{x, y, z}), velocidade
    (\texttt{vx, vy, vz}), aceleração (\texttt{ax, ay, az}) e atrito (\texttt{f}).
  \item \textbf{Visualização}: \texttt{s} (índice do tile no atlas) e \texttt{\_dead}
    (\emph{flag} de remoção).
  \item \textbf{Saúde}: \texttt{h} (\emph{hit points}), inicializado a 5 e
    sobrescrevível no \texttt{\_init}.
\end{itemize}

A interface comportamental é definida por seis métodos sobrescrevíveis:

\begin{description}
  \item[\texttt{\_init(param)}] \emph{Hook} de inicialização chamado pelo construtor.
    Permite à subclasse personalizar atributos sem reimplementar a cadeia de
    inicialização.
  \item[\texttt{\_update()}] Avança a física por \texttt{time\_elapsed} segundos. A
    implementação base aplica forças, integra a posição e testa colisão com paredes via
    \texttt{\_collides}.
  \item[\texttt{\_collides(x, z)}] Checa se a posição alvo está dentro de um tile
    sólido. Sobrescrito em \texttt{entity\_player\_t} para implementar a barreira do
    \emph{lock-and-clear}.
  \item[\texttt{\_check(other)}] Chamado para cada par de entidades cuja AABB se
    intersecta. Implementa \emph{pickup}, dano por contato, separação entre inimigos,
    entre outros.
  \item[\texttt{\_render()}] Emite primitivas para o \emph{renderer}
    (\texttt{push\_sprite}, \texttt{push\_light}).
  \item[\texttt{\_kill()}] Marca a entidade como \texttt{\_dead}, agendando sua remoção
    no fim do \emph{frame}.
\end{description}

\section{Subclasses de combate}

\subsection{\texttt{entity\_player\_t}}

O jogador é único em vários aspectos:

\begin{itemize}
  \item Sobrescreve \texttt{\_collides} para impedir saída de salas trancadas.
  \item Mantém estado de \emph{combo} (\texttt{combo}, \texttt{combo\_timer}) e
    pontuação.
  \item Gerencia o \emph{cooldown} do \emph{dash} (\texttt{\_dash\_cooldown},
    \texttt{\_dashing}) e da arma equipada (\texttt{\_last\_shot}).
  \item Em \texttt{\_update}, calcula o ângulo de mira a partir da posição do mouse e
    do \emph{offset} da câmera, e despacha o ataque pelo \texttt{type} da arma
    (\texttt{melee} ou \texttt{ranged}).
\end{itemize}

\subsection{Inimigos}

\begin{description}
  \item[\texttt{entity\_seeker\_t}] (HP 3). Persegue o jogador em linha reta com
    \emph{retarget} periódico. Comportamento mais simples, introduzido no Andar 1.
  \item[\texttt{entity\_shooter\_t}] (HP 5). Mantém distância de 24 unidades do
    jogador, e dispara em cone duas \texttt{entity\_enemy\_plasma\_t} ao detectá-lo.
    Aparece a partir do Andar 2.
  \item[\texttt{entity\_bomber\_t}] (HP 2). Persegue agressivamente; ao chegar a 20
    unidades de distância, ativa um \emph{fuse} de 0,8\,s e explode, causando dano em
    raio de 28 unidades. Introduzido no Andar 3.
  \item[\texttt{entity\_boss\_t}] (HP 70). Adversário do Andar 3, com duas fases:
    \begin{enumerate}
      \item Fase 1 (HP > 35): atira anéis de 7 plasmas em padrão rotativo.
      \item Fase 2 (HP $\leq$ 35): abandona o padrão rotativo, avança em \emph{melee}
        e invoca grupos de 3 \texttt{entity\_seeker\_t}.
    \end{enumerate}
\end{description}

\section{Subclasses de projétil e efeito}

\begin{description}
  \item[\texttt{entity\_plasma\_t}] Projétil do jogador (\emph{ranged}). Move-se em
    linha reta, hospeda \texttt{\_dmg} configurável pela arma, e é destruído ao colidir
    com parede ou inimigo.
  \item[\texttt{entity\_enemy\_plasma\_t}] Versão adversária, que só causa dano ao
    jogador.
  \item[\texttt{entity\_explosion\_t}] Efeito visual transitório de 1 segundo, com luz
    pulsante para \emph{feedback} de impacto.
  \item[\texttt{entity\_slash\_t}] \emph{Burst} de partículas com \emph{flash} de luz
    para o golpe \emph{melee}. \emph{Lifetime} curto, de 180\,ms.
  \item[\texttt{entity\_particle\_t}] Partícula com gravidade e \emph{bounce}, gerada
    por impactos.
\end{description}

\section{Subclasses de pickup e cenário}

\begin{description}
  \item[\texttt{entity\_health\_t}] \emph{Pickup} de poção (curativo de 1 HP). Adiciona
    ao \emph{slot} de consumível do jogador (capacidade 1). Spawn aleatório, com 50\%
    de chance por sala regular, ou garantido na sala do baú.
  \item[\texttt{entity\_weapon\_t}] \emph{Pickup} de arma. Sorteia um índice diferente
    da arma atualmente equipada e substitui ao toque.
  \item[\texttt{entity\_staircase\_t}] Escada de descida. Inicialmente inativa
    (\texttt{\_active = false}), é ativada por \texttt{count\_enemies} quando todos os
    inimigos da sala estão mortos. Ao ser tocada, dispara \texttt{show\_shop} (em
    andares não-finais) ou termina o jogo (no último andar).
\end{description}

\section{Lifecycle e remoção}

A remoção de entidades segue um \textit{deferred deletion} para preservar a
estabilidade da iteração:

\begin{enumerate}
  \item \texttt{\_kill()} marca \texttt{\_dead = true} e empurra a referência em
    \texttt{entities\_to\_kill}.
  \item Iterações subsequentes do \texttt{\_update}, \texttt{\_check} e
    \texttt{\_render} no mesmo \emph{frame} testam \texttt{\_dead} e pulam a entidade.
  \item No fim do \emph{frame}, \texttt{entities = entities.filter(\dots)} reconstrói
    o \emph{array} sem as entidades marcadas.
\end{enumerate}

\chapter{Máquina de Estados do Jogo}

\section{O \texttt{game\_state} global}

Uma única variável \texttt{game\_state} controla o comportamento do programa,
assumindo nove valores discretos. O \texttt{game\_tick} despacha lógica distinta para
cada estado; transições são disparadas por entrada do usuário, eventos da simulação
ou \texttt{setTimeout}s explícitos.

O \textbf{State Machine Diagram}, também conhecido como \emph{statechart}, comunica
simultaneamente os estados, transições e condições de guarda.

\diagfig{diagrams/state-machine.png}{Diagrama de Estados do \texttt{game\_state}}{state-machine}

\section{Catálogo de estados}

\begin{description}

  \item[\texttt{menu}] Tela de título. Exibe três opções (\textsc{JOGAR},
    \textsc{CRÉDITOS}, \textsc{SAIR}) navegáveis por setas ou WS, com confirmação por
    Enter ou Espaço. É o estado de entrada após o \emph{boot} e ponto de retorno
    garantido a partir de qualquer estado terminal.

  \item[\texttt{credits}] Tela de créditos com a equipe e referências de \emph{stack}.
    Implementa um \emph{grace period} de 300\,ms para ignorar o \emph{keyup} da
    própria tecla que abriu a tela. Qualquer tecla ou clique após esse intervalo
    retorna a \texttt{menu}.

  \item[\texttt{quit}] Tela de despedida ``SISTEMA DESLIGADO''. Tenta
    \texttt{window.close()}, que funciona no contexto Electron e falha silenciosamente
    no \emph{browser}. O \emph{fallback} para \texttt{menu} acontece via tecla ou
    clique.

  \item[\texttt{loading}] Estado transitório enquanto \texttt{regenerate\_atlas}
    aguarda o \texttt{Image} carregar. Acontece apenas no primeiro andar de cada
    sessão, pois chamadas subsequentes são síncronas com a imagem em cache.

  \item[\texttt{playing}] Estado central. O \texttt{game\_tick} executa as três fases
    do \emph{frame} (\texttt{update} $\rightarrow$ \texttt{collision} $\rightarrow$
    \texttt{render}), atualiza a HUD e checa o \emph{lock-and-clear}.

  \item[\texttt{paused}] Pausa por \textsc{ESC}. Sobrescreve os \emph{handlers} de
    teclado para ignorar tudo exceto novo \textsc{ESC}. O \texttt{game\_tick} faz
    \emph{early return} sem processar entidades nem renderizar; o \texttt{canvas}
    mantém o último \emph{frame} desenhado, e o \texttt{hud\_ctx} recebe uma camada
    semi-transparente com o texto ``PAUSADO''.

  \item[\texttt{shop}] Loja entre andares, ativada quando o jogador toca a escada em
    um andar não-final. Apresenta três \emph{upgrades} sorteados de um \emph{pool}
    persistente de cinco. \emph{Upgrades} comprados são removidos do \emph{pool} para
    o resto da \emph{run}.

  \item[\texttt{gameover}] Tela de derrota. Disparada pelo \texttt{setTimeout} de
    1500\,ms agendado em \texttt{entity\_player\_t.\_kill}. Suporta dois caminhos:
    \textsc{Enter} reinicia a \emph{run}, e \textsc{Esc} retorna ao menu.

  \item[\texttt{victory}] Tela de vitória. Disparada pelo \texttt{setTimeout} de
    2000\,ms em \texttt{entity\_boss\_t.\_kill}. Mostra estatísticas finais (tempo,
    kills, score, arma) e retorna ao menu por qualquer entrada.

\end{description}

\section{Transições críticas}

\subsection{\emph{Setup} de \emph{handlers} por estado}

Cada estado instala seu próprio conjunto de \emph{handlers}: \texttt{onkeydown},
\texttt{onkeyup}, \texttt{onmousedown} e \texttt{onmouseup}. A função
\texttt{bind\_input} restaura os \emph{handlers} de \emph{gameplay} ao entrar em
\texttt{playing}. Outros estados usam o \texttt{\_bind\_dismiss}, um \emph{helper}
que adiciona o \emph{grace period} e responde em \texttt{keyup} ou \texttt{mouseup},
ou implementam \emph{handlers} ad-hoc.

\subsection{Race conditions com \texttt{setTimeout}}

A morte do jogador agenda \texttt{show\_gameover\_screen} em 1500\,ms, e a morte do
\emph{boss} agenda \texttt{game\_victory} em 2000\,ms. Se as duas mortes acontecerem
no mesmo \emph{frame} (cenário plausível, em que a explosão final do \emph{boss}
mata o jogador), os dois \texttt{setTimeout}s ficam pendentes e disparam em janelas
diferentes. As transições estão protegidas por guardas idempotentes:

\begin{itemize}
  \item \texttt{show\_gameover\_screen} ignora se \texttt{game\_state} já é
    \texttt{'gameover'} ou \texttt{'victory'}.
  \item \texttt{game\_victory} ignora se \texttt{game\_state} já é \texttt{'victory'}.
\end{itemize}

A vitória, porém, \emph{sobrescreve} a derrota. A interpretação de design é direta:
se o jogador matou o \emph{boss}, o objetivo foi cumprido, independentemente de uma
morte contemporânea.

\subsection{Reset de \texttt{keys[]}}

Toda transição que entra em \texttt{playing} (a partir de \texttt{loading},
\texttt{shop} ou \texttt{gameover}) chama
\texttt{for (var k in keys) keys[k] = 0;} antes de devolver o controle. A linha evita
o problema de \emph{stuck keys}, em que uma tecla pressionada durante o estado
anterior ficaria registrada como ativa quando o jogo retomasse.

\chapter{Game Loop}

O \emph{game loop} do Void Descent opera em dois níveis: \textbf{micro} (dentro de uma
sala, segundo a segundo) e \textbf{macro} (ao longo de uma \emph{run}, minuto a
minuto). \textbf{Activity Diagrams} representam fluxos com decisões e iteração, e são
o tipo UML adequado para ambos.

\section{Loop Micro: lock-and-clear}

O \emph{loop} micro é o ciclo \emph{lock-and-clear} clássico de roguelikes. Ao entrar
em uma sala com inimigos, as ``portas'' trancam, implementadas como restrição lógica
de movimento do jogador, não como tiles físicos. O combate continua até que todos os
inimigos \emph{originalmente associados àquela sala} sejam eliminados. A associação é
feita no \emph{spawn} via \texttt{e.\_room\_idx = i}, evitando o \emph{exploit}
clássico em que um inimigo escapando para o corredor liberaria a sala falsamente.

A duração típica de um ciclo micro é de 5 a 10 segundos. As ações disponíveis ao
jogador acontecem em paralelo:

\begin{itemize}
  \item \textbf{Movimento}: WASD, $\sim$128 unidades por segundo, escalado por
    \texttt{player\_speed\_mult}.
  \item \textbf{Ataque}: o clique esquerdo dispara projétil (\emph{ranged}) ou aplica
    dano em arco (\emph{melee}).
  \item \textbf{Dash}: \textsc{Espaço} ativa 0,2\,s de invencibilidade com aceleração
    de 400 unidades por segundo, com \emph{cooldown} de 1,5\,s.
  \item \textbf{Cura}: \textsc{Q} consome 1 poção e restaura 1 HP, se houver
    \emph{slot} cheio e o jogador estiver abaixo do HP máximo.
\end{itemize}

\diagfig{diagrams/micro-loop.png}{Diagrama de Atividade do ciclo de combate por sala}{micro-loop}

\section{Loop Macro: progressão da run}

O \emph{loop} macro segue a curva de dificuldade de três atos prescrita no GDD:

\begin{enumerate}
  \item \textbf{Andar 1, Antecâmara} (paleta ciano). Apresentação suave: apenas
    \texttt{Seekers}, salas espaçosas, drops generosos. Ensino implícito da mecânica
    básica de movimentar-se e atacar.

  \item \textbf{Andar 2, Subsistemas} (paleta amarela). Introdução dos
    \texttt{Shooters}. Pela primeira vez, o jogador precisa usar o \emph{dash}
    ofensivamente para desviar de projéteis. A complexidade tática aumenta, sem
    mudança fundamental nas ferramentas.

  \item \textbf{Andar 3, Núcleo} (paleta magenta). Clímax: \texttt{Bombers}
    introduzidos, arena curta, culminando no boss em duas fases. A brevidade é
    proposital, pois após a extensão do Andar 2 o ritmo acelera deliberadamente em
    direção ao confronto final.
\end{enumerate}

Entre andares (após Andar 1 e Andar 2), o jogador acessa a \emph{shop} de
\emph{upgrades}. O \emph{pool} de cinco \emph{upgrades} se esvazia ao longo da
\emph{run}; combinado ao \emph{permadeath}, isso produz escolhas únicas que não se
repetem.

\diagfig{diagrams/macro-loop.png}{Diagrama de Atividade da progressão completa da \emph{run}}{macro-loop}

\section{Sequência interna de um \emph{frame}}

A função \texttt{game\_tick}, invocada por \texttt{requestAnimationFrame} a
$\sim$60\,Hz, executa três fases estritamente ordenadas:

\begin{enumerate}
  \item \textbf{Update}: para cada entidade viva, \texttt{\_update} atualiza posição e
    estado interno. A colisão com paredes acontece dentro de \texttt{\_update}, via
    \texttt{\_collides}. Esta fase é puramente local, sem interação entre entidades.

  \item \textbf{Collision}: \texttt{check\_collisions\_grid} executa o \emph{broad
    phase} via \emph{spatial hash}. Cada entidade é \emph{bucketizada} numa célula de
    16$\times$16 unidades. Para cada par de entidades em células vizinhas (3$\times$3
    \emph{neighborhood}), o teste AABB é executado e, em caso de interseção,
    \texttt{e1.\_check(e2)} e \texttt{e2.\_check(e1)} são invocados. A deduplicação de
    pares é feita por um campo temporário \texttt{\_gid}, atribuído pelo índice no
    \emph{array} durante a construção do \emph{grid}.

  \item \textbf{Render}: cada entidade emite suas primitivas via \texttt{push\_sprite}
    e \texttt{push\_light}. Após o \emph{loop}, \texttt{renderer\_end\_frame} executa
    o \emph{draw call} único do WebGL.
\end{enumerate}

A separação em fases foi introduzida durante a refatoração do \emph{spatial grid}. O
modelo anterior, um \emph{loop} aninhado $O(n^2)$ que intercalava \texttt{\_update},
\texttt{\_check} e \texttt{\_render} em cada iteração, não escalava para os
$\sim$400 inimigos por \emph{run} do modo de dificuldade alta.

Após as três fases, a HUD é desenhada no \texttt{hud\_ctx} (\emph{canvas} 2D
separado), \texttt{check\_room\_lock} verifica transições de sala, e o
\emph{deferred deletion} remove as entidades marcadas mortas.

\diagfig{diagrams/frame-sequence.png}{Diagrama de Sequência da execução interna de \texttt{game\_tick} no estado \texttt{playing}}{frame-sequence}

\chapter{Sistemas Procedurais}

A premissa fundamental do Void Descent é que \emph{nenhum asset é importado}. Texturas,
geometria, áudio, comportamento de inimigos e mapas são gerados algoritmicamente a cada
\emph{run}. Este capítulo descreve os três sistemas procedurais centrais usando
\textbf{Sequence Diagrams}, que representam interação ordenada entre componentes ao
longo do tempo, o tipo UML adequado para descrever \emph{processos} multi-etapas.

\section{Geração de \emph{dungeon}}

\diagfig{diagrams/dungeon-gen.png}{Diagrama de Sequência da geração procedural de um andar}{dungeon-gen}

\subsection{Algoritmo de \emph{room placement}}

A função \texttt{dungeon\_generate(floor\_id)} segue cinco etapas:

\begin{enumerate}
  \item \textbf{Seeding}. \texttt{random\_seed(0xBADC0DE1 + floor\_id)} reinicia o LCG.
    Andares com mesmo \emph{seed} produzem o mesmo mapa, garantindo determinismo do
    \emph{layout} dentro de uma \emph{run}. Isso é importante para que
    \texttt{count\_enemies\_in\_room} identifique consistentemente a mesma sala em
    \emph{frames} diferentes.

  \item \textbf{Posicionamento de salas}. Em um \emph{loop} com até 500 tentativas,
    são sorteados retângulos com $w, h \in [7, 11]$ tiles e $x, y$ válidos no mapa
    $64 \times 64$. Cada candidato é rejeitado se sobrepõe (com \emph{gap} mínimo de
    1 tile) qualquer sala já aceita. O número alvo é $4 + \mathit{floor\_id}$ salas,
    exceto no Andar 3, que tem 5.

  \item \textbf{Tipagem}. \texttt{rooms[0]} é a sala inicial; a última é
    \texttt{chest} (Andar 1 e 2) ou \texttt{boss} (Andar 3). No Andar 3 a segunda
    sala também é \texttt{chest}, garantindo equipamento antes do confronto final.

  \item \textbf{Carving} dos corredores. Para cada par de salas adjacentes na lista,
    um corredor em forma de L é escavado conectando os centroides, com escolha
    aleatória da ordem (horizontal-primeiro ou vertical-primeiro).

  \item \textbf{Construção da geometria}. Cada tile é convertido em chamada a
    \texttt{push\_floor} (chão) ou \texttt{push\_block} (parede). O resultado é
    \emph{baked} no \emph{vertex buffer} do WebGL. Não há atualização incremental
    durante o \emph{gameplay}, o que permite que \texttt{level\_num\_verts} marque
    o fim da geometria estática.
\end{enumerate}

\subsection{Validação}

O posicionamento por descarte tem chance teórica de gerar mapas onde uma sala fique
ilhada. Na prática, com 500 tentativas e o algoritmo de conexão sequencial
\texttt{i $\to$ i+1}, todas as salas ficam alcançáveis a partir da inicial via grafo
linear de corredores.

\section{Síntese de áudio}

\diagfig{diagrams/audio-synth.png}{Diagrama de Sequência do \emph{pipeline} de síntese de áudio}{audio-synth}

\subsection{Pipeline}

A \texttt{audio\_init} é assíncrona:

\begin{enumerate}
  \item Cria o \texttt{AudioContext}, inicialmente \emph{suspended} pela política de
    \emph{autoplay}, e resumido na primeira interação do usuário.
  \item Cria \texttt{audio\_master\_gain} (um \texttt{GainNode}) e o conecta ao
    \texttt{destination}. Todas as fontes posteriores passam por este nó,
    permitindo \emph{mute} unificado via
    \texttt{audio\_master\_gain.gain.value = 0}.
  \item Dispara a síntese assíncrona da música. \texttt{sonantxr\_generate\_song}
    processa cada \emph{track} em \emph{chunks} de 33\,ms, intercalando com
    \texttt{setTimeout(0)} para não bloquear a \emph{main thread} excessivamente.
  \item Em paralelo, dispara a síntese de oito SFX (\emph{shoot}, \emph{hit},
    \emph{hurt}, \emph{beep}, \emph{pickup}, \emph{terminal}, \emph{explode},
    \emph{dash}). Cada um é síncrono, mas curto.
\end{enumerate}

\subsection{Estrutura interna do sintetizador}

Cada nota processada em \texttt{SoundGenerator.\_genSound} executa as seguintes
etapas, sobre os $\sim$10\,000 \emph{samples} de duração da nota:

\begin{itemize}
  \item Dois osciladores (\texttt{osc1}, \texttt{osc2}) com formas selecionáveis
    (\textit{sin}, \textit{square}, \textit{saw}, \textit{tri}) e \emph{detune}
    aplicável.
  \item Envelope ADSR (\texttt{attack}, \texttt{sustain}, \texttt{release})
    calculado por \emph{sample}.
  \item Filtro biquad de estado configurável (\textit{lopass}, \textit{hipass},
    \textit{bandpass}, \textit{notch}) com ressonância.
  \item LFO modulando a frequência do filtro.
  \item \emph{Pan} e \emph{delay} aplicados em pós-processamento sobre o
    \texttt{Float32Array} resultante.
\end{itemize}

A faixa ambiente atual usa um único \emph{track} de \emph{drone} com
$rowLen = 11025$ e $songLen = 16$\,s, gerando um \emph{loop} sem cortes audíveis.

\section{Atlas de tiles e \emph{tinting}}

O atlas de tiles é embarcado no código como \emph{string} \texttt{base64}
(\texttt{TILES\_B64} em \texttt{game.js}). No \emph{boot} de cada andar,
\texttt{tint\_atlas} aplica um \emph{hue shift} ao atlas inteiro pixel-a-pixel,
produzindo a paleta característica de cada andar: ciano para o Andar 1
($\Delta H = 0{,}55$), laranja para o Andar 2 ($\Delta H = 0{,}15$) e magenta
para o Andar 3 ($\Delta H = 0{,}75$). O resultado é uma textura
$\mathit{tile\_atlas\_size}^2$ vinculada ao GL via \texttt{renderer\_bind\_image}.

A conversão de espaço é \texttt{RGB} $\rightarrow$ \texttt{HSL} $\rightarrow$
\texttt{shifted HSL} $\rightarrow$ \texttt{RGB}, com saturação amplificada em 40\%
para preservar o \emph{look} \emph{neon} sobre fundo preto.

\chapter{Build e Deploy}

\section{Pipeline de empacotamento}

O Void Descent é distribuído como executável \texttt{.exe} portátil para Windows
x64, empacotado via Electron e \texttt{electron-builder}. Para representar a
topologia de \emph{nodos} físicos e lógicos, e os \emph{artifacts} que circulam
entre eles, o tipo de diagrama UML adequado é o \textbf{Deployment Diagram}.

\diagfig{diagrams/deployment.png}{Diagrama de Deployment do \emph{pipeline} de \emph{build} e \emph{runtime}}{deployment}

\section{Comandos de \emph{build}}

\subsection{Desenvolvimento local}

Para rodar a versão de desenvolvimento, com \emph{hot reload} manual via
\emph{refresh} do Electron:

\begin{verbatim}
npm install
npm start
\end{verbatim}

\subsection{Build do \texttt{.exe} para Windows}

A configuração em \texttt{package.json} define o \emph{target} \texttt{portable},
um executável auto-extraível sem instalador:

\begin{verbatim}
npm run build:win
\end{verbatim}

A saída é \texttt{dist/VoidDescent-1.0.0.exe}, com $\sim$73\,MB. Em paralelo,
\texttt{electron-builder} produz \texttt{dist/win-unpacked/} (a árvore de arquivos
antes do empacotamento, útil para \emph{debug}) e arquivos auxiliares
(\texttt{builder-debug.yml}, \emph{blockmaps}) que são ignorados pelo
\texttt{.gitignore}.

\section{Cross-compile do macOS}

A geração do \texttt{.exe} para Windows acontece no \emph{host} macOS, sem
necessidade de Wine ou máquina virtual Windows. O \emph{target} \texttt{portable} é
uma das poucas configurações que o \texttt{electron-builder} cross-compila
nativamente; alvos NSIS (instalador) ou MSI exigem Wine no \texttt{PATH}.

\section{Por que 73\,MB?}

A composição aproximada do binário \texttt{.exe} é a seguinte:

\begin{tabular}{lr}
  \toprule
  \textbf{Componente} & \textbf{Tamanho} \\
  \midrule
  Chromium (V8 e Blink) & $\sim$60\,MB \\
  Node.js \emph{runtime} & $\sim$10\,MB \\
  Locales (ICU) & $\sim$5\,MB \\
  Glue do Electron & $\sim$3\,MB \\
  Bundle do jogo (\texttt{app.asar}) & $\sim$80\,KB \\
  \bottomrule
\end{tabular}

O ``jogo em si'' ocupa $0{,}11\%$ do binário. Os outros $99{,}89\%$ são o
\emph{runtime} necessário para hospedar uma página WebGL como aplicação
\emph{desktop}. Esta foi a escolha explícita do GDD; a alternativa seria usar o
\emph{webview} do sistema (Tauri ou Neutralino reduziriam o tamanho para 2 a
10\,MB), mas o requisito de portabilidade absoluta e reprodutibilidade pesou na
decisão por Electron.

\section{Distribuição}

O \texttt{.exe} é \emph{não-assinado}, pois não há certificado de \emph{code
signing} de desenvolvedor Microsoft. Na primeira execução em Windows 10 ou 11, o
\emph{SmartScreen} exibe o aviso ``O Windows protegeu seu PC''. A interação
esperada do usuário é:

\begin{enumerate}
  \item Clicar em \textit{Mais informações}.
  \item Clicar em \textit{Executar mesmo assim}.
\end{enumerate}

Após essa primeira execução, o \emph{hash} do executável é registrado localmente e
o aviso não reaparece. O modelo é coerente com a distribuição via \emph{releases}
do GitHub, o canal escolhido para entrega.

\section{Repositório}

O código-fonte e o \texttt{.exe} compilado estão disponíveis em:

\begin{center}
  \url{https://github.com/luisg0c/Void-Descent}
\end{center}

A história de \emph{commits} reflete a progressão real do projeto entre 03 de
março e 09 de maio de 2026, com mensagens em formato \emph{Conventional Commits}.

\chapter{Backlog Técnico, Decisões e Aprendizados}

Este capítulo documenta o caminho real do projeto: as decisões que tomamos, as
ferramentas que escolhemos, o que ficou de fora e o que aprendemos no processo. Nem
tudo era óbvio no GDD original; algumas escolhas só fizeram sentido ao bater em
problemas concretos.

\section{Decisões técnicas}

\subsection{JavaScript puro sem \emph{framework}}

Considerações descartadas: Phaser, PixiJS, Three.js. Optamos por JavaScript puro com
WebGL direto pelos seguintes motivos:

\begin{itemize}
  \item \textbf{Tamanho do bundle}. Phaser sozinho tem $\sim$1\,MB, e PixiJS tem
    $\sim$500\,KB. Nosso código fonte completo tem 77\,KB. Um \emph{framework}
    representaria mais de 90\% do peso, contradizendo a proposta de procedural
    minimalista.
  \item \textbf{Aprendizado}. Implementar o \emph{pipeline} de \emph{shaders}, o
    \emph{tinting} do atlas e a iluminação por vértice à mão deu domínio sobre o que
    cada parte faz.
  \item \textbf{Sem cerimônia de \emph{boot}}. \emph{Frameworks} têm \emph{loaders},
    cenas e sistema de eventos próprio. Para um jogo de tela única e fluxo linear, é
    \emph{overkill}.
\end{itemize}

\subsection{Atlas de tiles embarcado em \emph{base64}}

A \texttt{TILES\_B64} em \texttt{game.js} é uma \emph{string} \emph{base64} de
$\sim$3\,KB, que decodifica para um PNG 1024$\times$16. Considerações descartadas:

\begin{itemize}
  \item \textbf{PNG externo}. Adicionaria uma dependência de carregamento assíncrono
    e quebraria o \emph{single-file deploy}. O jogo roda abrindo \texttt{index.html}
    sem servidor.
  \item \textbf{Geração 100\% procedural via \emph{shaders}}. Investigamos, mas o
    custo de tempo de implementação superava o benefício. O atlas pintado à mão é
    \emph{tile-based} e mantém a estética \emph{neon} desejada.
\end{itemize}

A escolha pelo atlas embarcado preserva a estética \emph{neon} sem comprometer a
\emph{distribution story} de \emph{single-page}.

\subsection{Sonant-X em vez de Web Audio API direta}

A Web Audio API permite síntese arbitrária via \texttt{OscillatorNode},
\texttt{BiquadFilterNode} e outros nós. Implementar oito SFX e uma faixa ambiente
diretamente sobre essa API resultaria em centenas de linhas de \emph{boilerplate}.
O Sonant-X é uma biblioteca de $\sim$200 linhas que abstrai oscilador, envelope,
filtro, LFO e \emph{delay} num único \texttt{instr} declarativo. Encurta o tempo de
iteração das peças sonoras de horas para minutos.

\subsection{Electron em vez de Tauri ou Neutralino}

O GDD especifica Electron, e a implementação seguiu essa direção. Avaliamos
alternativas durante o \emph{packaging}:

\begin{tabular}{lrl}
  \toprule
  \textbf{Tecnologia} & \textbf{Tamanho .exe} & \textbf{Custo de migração} \\
  \midrule
  Electron 32 (escolhido) & 73\,MB & 0 \\
  Tauri 2 (Rust webview) & $\sim$5--10\,MB & Reescrita do \emph{boot}, Rust toolchain \\
  Neutralino & $\sim$2\,MB & Reescrita do \emph{boot}, depende de WebView2 \\
  Native C++ + OpenGL & $<$500\,KB & Reescrita completa \\
  \bottomrule
\end{tabular}

A redução de Electron para Neutralino encolheria o \texttt{.exe} em $\sim$35
vezes, mas o prazo próximo e o risco de \emph{webview} ausente em máquinas antigas
pesaram contra a migração.

\subsection{\emph{Spatial grid} 32$\times$32 com células de 16}

Decisão tomada após o modo de dificuldade alta elevar a contagem de inimigos para
$\sim$400 por \emph{run}. O \emph{loop} aninhado $O(n^2)$ original processava
80.000 pares por \emph{frame}, o suficiente para causar quedas de 60\,Hz para
$\sim$30\,Hz em telas com muitos inimigos.

A escolha de células de 16 unidades é informada pelo tamanho da AABB da entidade
(9 unidades). Cada entidade ocupa, no máximo, duas células, e a busca de vizinhos
em 3$\times$3 captura todas as colisões possíveis. Células menores aumentariam o
\emph{overhead} de \emph{hashing}; maiores diluiriam o ganho.

\subsection{\emph{Lock-and-clear} via colisão lógica em vez de \emph{tiles} físicos}

A leitura natural do GDD, com a expressão ``portas trancam'', sugeria adicionar
\emph{tiles} de parede dinamicamente nos corredores ao entrar em sala com
inimigos. Descartado:

\begin{itemize}
  \item O \emph{vertex buffer} do WebGL é compilado uma vez no \texttt{load\_level}
    e permanece estático durante o \emph{gameplay}. Modificar tiles requer
    recompilar a geometria, custoso.
  \item Solução adotada: a função \texttt{entity\_player\_t.\_collides} é
    sobrescrita para checar se a posição alvo está dentro do retângulo da sala
    atualmente trancada. Sem custo de geometria, sem mudança visual.
\end{itemize}

A solução tem um \emph{trade-off}: barreiras invisíveis. O efeito é mitigado com
\emph{feedback} sonoro e visual via \texttt{terminal\_show\_notice}.

\section{Backlog: features descartadas ou postergadas}

\begin{description}

  \item[Tutorial explícito] \emph{Fora de escopo}. O GDD prescreve
    \emph{teaching through play}. A primeira sala do Andar 1 ensina movimentação
    contra \texttt{Seekers} sem nenhuma tela de instrução, e a HUD é o único
    material didático.

  \item[Sistema de \emph{achievements}] \emph{Descartado}. Adicionaria persistência
    entre \emph{runs}, contradizendo o \emph{permadeath} estrito e exigindo
    armazenamento local que o GDD evita.

  \item[Multiplayer cooperativo] \emph{Fora de escopo}. Não previsto no GDD.

  \item[Mais de três andares] \emph{Limitado pelo GDD}. A estrutura de três atos
    (apresentação, teste, clímax) é canônica no documento e foi mantida.

  \item[Validação \emph{flood-fill} do mapa] \emph{Implementado parcialmente}. O
    GDD prevê uma fase de validação por \emph{flood-fill} para garantir que todas
    as salas sejam alcançáveis. Na prática, o algoritmo de conexão sequencial entre
    salas (\texttt{i $\to$ i+1}) já garante alcançabilidade por construção, e a
    validação se mostrou redundante. Documentado como TODO se geradores futuros
    forem mais sofisticados, como Voronoi.

  \item[Salvar e carregar \emph{run}] \emph{Explicitamente fora}. \emph{Permadeath}
    é decisão de design central; persistência quebraria o ciclo.

  \item[Inimigos novos no Andar 3] \emph{Implementado}. \texttt{Bombers} entram
    apenas no Andar 3, conforme o GDD.

  \item[\emph{Boss} com mais de duas fases] \emph{Considerado, descartado}. Duas
    fases é o ponto de equilíbrio observado em \emph{playtesting}: pune
    complacência sem exaurir o jogador. Três fases prolongariam o confronto e
    diluiriam o impacto.

\end{description}

\section{Ferramentas e \emph{assets}}

\subsection{Estética visual}

\begin{itemize}
  \item \textbf{Atlas de tiles}: um único PNG 1024$\times$16, embarcado como
    \emph{base64} em \texttt{game.js}. Os 16 tiles cobrem chão, parede, inimigos,
    armas, explosões e ícones de UI. Foi o único \emph{asset} pintado à mão.
  \item \textbf{Paletas por andar}: aplicadas via \emph{hue shift} algorítmico
    (HSL) sobre o atlas no \texttt{regenerate\_atlas}. Não há atlas separado por
    andar; todos derivam do mesmo PNG.
  \item \textbf{Sprites de inimigos}: 4 \emph{frames} para \texttt{seeker}
    (animação) e sprite estático para \texttt{shooter}, \texttt{bomber} e
    \texttt{boss}.
\end{itemize}

\subsection{Áudio}

\begin{itemize}
  \item \textbf{Sintetizador}: Sonant-X (\emph{port} JavaScript de Sonant Live).
  \item \textbf{Faixa ambiente}: \emph{drone} em Em, com 16\,s de \emph{loop}
    perfeito. Iterada cinco vezes até atingir uma textura aceitável; a primeira
    versão era \emph{rock metal} agressivo, descartado por conflitar com a estética
    de terminal sci-fi.
  \item \textbf{SFX}: 8 \emph{patches} (\emph{shoot}, \emph{hit}, \emph{hurt},
    \emph{beep}, \emph{pickup}, \emph{terminal}, \emph{explode}, \emph{dash}).
\end{itemize}

\subsection{Cadeia de \emph{tools}}

\begin{itemize}
  \item \textbf{IDE}: VS Code.
  \item \textbf{Servidor de desenvolvimento}: \texttt{python3 -m http.server} no
    \texttt{macOS}, para servir \texttt{index.html} durante iteração.
  \item \textbf{Empacotamento}: \texttt{electron-builder} 25 com \emph{target}
    \texttt{portable}.
  \item \textbf{Documentação}: \LaTeX{} (\texttt{latexmk}) e PlantUML para os
    diagramas UML. \emph{Pipeline}: \texttt{.puml} $\to$ \texttt{.png} $\to$
    \texttt{\textbackslash includegraphics} $\to$ PDF final.
  \item \textbf{Versionamento}: Git e GitHub. Histórico de \emph{commits} no
    formato \emph{Conventional Commits}.
\end{itemize}

\section{Aprendizados técnicos}

\subsection{\emph{Race conditions} em \texttt{setTimeout} aninhados}

A morte do jogador agenda \texttt{show\_gameover\_screen} em 1500\,ms, e a morte
do \emph{boss} agenda \texttt{game\_victory} em 2000\,ms. Se ambas as mortes
acontecem no mesmo \emph{frame}, os dois \texttt{setTimeout}s ficam pendentes e
podem disparar em qualquer ordem. Aprendizado: toda transição de estado precisa
ser \emph{idempotente} ou guardada por \emph{state checks}; não confiar na ordem
de execução de \texttt{setTimeout}.

\subsection{Coordenadas do \emph{mouse} requerem \texttt{getBoundingClientRect}}

A primeira implementação usava \texttt{ev.clientX / canvas.clientWidth *
canvas.width}, ignorando o \emph{offset} do \emph{canvas} no \emph{viewport}.
Funcionava quando o \emph{canvas} ocupava 100\% da janela, mas falhava quando
havia margens (típico em \emph{flexbox} centralizado). A correção via
\texttt{canvas.getBoundingClientRect()} foi descoberta empiricamente, depois que
o jogador precisava apontar para fora da arma para mirar.

\subsection{\texttt{textAlign} do \emph{canvas} 2D persiste entre chamadas}

Bug clássico: setamos \texttt{textAlign = 'right'} para o \texttt{POCAO} e
esquecemos de resetar para \texttt{'left'} antes de desenhar
\texttt{DASH: PRONTO}, que renderizava fora da tela. Aprendizado: estados de
contexto 2D (cor, alinhamento, fonte) são \emph{sticky}; sempre setar
explicitamente o que cada bloco precisa, ou usar
\texttt{ctx.save()/restore()}.

\subsection{\emph{Lock-and-clear} precisa de \emph{tracking} explícito por entidade}

Versão ingênua: contar inimigos cuja posição está dentro do retângulo da sala.
Falha: um \texttt{seeker} perseguindo o jogador pode sair temporariamente da sala
pelo corredor, e a contagem cai para zero, liberando a sala falsamente. Solução:
tagar cada inimigo no \emph{spawn} com \texttt{e.\_room\_idx = i} e contar por
tag, não por posição. Aprendizado: para invariantes de \emph{game state},
prefira \emph{tags} persistentes a \emph{queries} espaciais.

\subsection{\emph{Loop} de música ambiente exige alinhamento de \emph{samples}}

A primeira versão da música tinha um pequeno \emph{gap} a cada \emph{loop}
($\sim$0,8\,s) porque o \texttt{songLen} (em segundos) não era múltiplo exato do
\texttt{rowLen} (em \emph{samples}). Solução: escolher \texttt{rowLen} e o
número de \emph{patterns} de modo que
$\mathit{patterns} \times 32 \times \mathit{rowLen} =
\mathit{songLen} \times 44100$ exatamente. No nosso caso,
$2 \times 32 \times 11025 = 16 \times 44100 = 705{.}600$ \emph{samples}.

\subsection{Cross-compile macOS para Windows funciona, com ressalvas}

\texttt{electron-builder} cross-compila do \texttt{macOS} para Windows sem Wine
\emph{somente} se o \emph{target} for \texttt{portable}. Alvos NSIS, MSI ou AppX
exigem Wine ou execução em VM Windows. Aprendizado: a escolha do \emph{target} é
tão crítica quanto a do \emph{framework} de empacotamento.

\section{Mudanças de escopo em relação ao GDD}

Durante o desenvolvimento, algumas especificações do GDD foram reinterpretadas,
estendidas ou explicitamente divergiram. Esta seção lista cada caso, com a
justificativa.

\subsection{Divergências de \emph{gameplay}}

\subsubsection*{Velocidade dos \texttt{Seekers}}

\textbf{GDD}: ``\texttt{Seekers} são inimigos \emph{melee} lentos cujo único
comportamento é caminhar na direção do jogador em linha reta.''\\
\textbf{Implementação}: aceleração de 210 unidades por segundo, agressivos.\\
\textbf{Razão}: durante o \emph{playtesting}, foi adotado um modo de dificuldade
alta a pedido explícito da equipe. Manter \texttt{Seekers} lentos nesse modo
tornaria o Andar 1 trivial. A versão original (lentos) seria adequada para um
modo ``normal'', mas não foi separada em uma \emph{flag}.

\subsubsection*{Comportamento do \emph{boss} na fase 2}

\textbf{GDD}: ``Na fase 2, ele abandona os projéteis, invoca \texttt{Seekers} como
reforço para pressionar o jogador e avança diretamente em \emph{melee}.''\\
\textbf{Implementação}: na fase 2, o \emph{boss} mantém um disparo de plasma único
a cada 0,7\,s, avança em \emph{melee} e invoca pares de \texttt{Seekers} a cada
3,0\,s.\\
\textbf{Razão}: sem qualquer projétil, a fase 2 ficaria mais fácil que a fase 1,
contra-intuitivo. Mantemos um disparo unitário para preservar a pressão
sem saturar a arena.

\subsubsection*{Densidade de inimigos por sala}

\textbf{GDD}: não especifica número exato; prescreve aumento progressivo entre
andares.\\
\textbf{Implementação}: 10 a 13 inimigos por sala no Andar 1, 14 a 17 no Andar 2
e 18 a 21 no Andar 3.\\
\textbf{Razão}: a contagem original (2 a 7 por sala) tornava o jogo trivial em
\emph{playtesting}; o aumento foi calibrado iterativamente.

\subsubsection*{Drops generosos no Andar 1, escassos no Andar 2}

\textbf{GDD}: ``Os drops nesse andar [Andar 1] são generosos. \dots\ no Andar 2,
a generosidade diminui de forma perceptível.''\\
\textbf{Implementação}: taxa uniforme de 50\% de chance de poção por sala em
todos os andares, com poção garantida na sala do baú.\\
\textbf{Razão}: simplificação do balanço. A diferenciação por andar adicionaria
complexidade sem ganho proporcional dado o escopo do P2.

\subsubsection*{Sala do \emph{boss} no Andar 3}

\textbf{GDD}: ``[Andar 3] possui apenas 3 salas regulares, incluindo uma sala de
baú, seguidas pela arena final do \emph{boss}.''\\
\textbf{Implementação}: 4 salas regulares mais a arena do \emph{boss}, totalizando
5 salas no Andar 3.\\
\textbf{Razão}: ajuste empírico durante \emph{playtesting}. Com 3 salas, o Andar
3 passava rápido demais e a expansão da dificuldade não tinha tempo de se sentir.

\subsection{Divergências de \emph{interface}}

\subsubsection*{Tecla \texttt{E} para interagir}

\textbf{GDD}: ``\texttt{E}: Interagir (coletar item ou usar escada).''\\
\textbf{Implementação}: \emph{auto-pickup} ao tocar; a escada também ativa ao
toque.\\
\textbf{Razão}: decisão consciente de UX. Em \emph{runs} curtas com muitos
inimigos, a fricção de uma tecla extra para coletar prejudicaria o \emph{flow
state}. \texttt{E} ficou disponível como \emph{slot} reservado para futuras
ferramentas de interação, como portas opcionais ou NPCs.

\subsubsection*{Sistema de pontuação}

\textbf{GDD}: $\mathit{score} = \mathit{kills} \times \mathit{combo} +
\mathit{itens} \times 100$.\\
\textbf{Implementação}: \texttt{score += base\_score $\times$ combo} a cada
\emph{kill} (incremental, não baseado em \emph{kills} totais), com bônus de
$+100$ por \emph{pickup}.\\
\textbf{Razão}: a fórmula incremental dá \emph{feedback} imediato no HUD a cada
abate, em vez de calcular ao final. É semanticamente equivalente, pois
\emph{kills} altos com combo alto produzem score alto. A versão do GDD exigiria
recálculo a cada \emph{kill} ou ao final, sem benefício de UX.

\subsubsection*{\emph{Slot} de \emph{passivo} dedicado}

\textbf{GDD}: ``1 \emph{slot} de \emph{passivo} \dots\ ocupado por modificadores
que persistem durante toda a \emph{run} (como velocidade aumentada, dano extra
ou HP adicional).''\\
\textbf{Implementação}: o sistema de \emph{shop} entre andares cumpre essa
função. \texttt{REPARO}, \texttt{OVERCLOCK}, \texttt{AMPLIFIER},
\texttt{ACELERADOR} e \texttt{CADÊNCIA} são todos modificadores persistentes.\\
\textbf{Razão}: a estrutura é equivalente ao GDD em termos funcionais, mas
empacotada como \emph{shop} (escolha ativa entre opções) em vez de \emph{drop}
aleatório (escolha passiva). \emph{Trade-off}: ganha agência do jogador, perde
descoberta orgânica.

\subsection{Divergências de \emph{geração}}

\subsubsection*{Validação por \emph{flood-fill}}

\textbf{GDD}: ``Após essa etapa, uma validação por \emph{flood-fill} percorre o
mapa a partir do ponto de \emph{spawn} do jogador para garantir que todas as
salas sejam alcançáveis.''\\
\textbf{Implementação}: validação não implementada; alcançabilidade garantida
por construção (o algoritmo conecta sala $i$ a $i+1$ sequencialmente).\\
\textbf{Razão}: redundância. A construção sequencial impede salas isoladas. Se
o gerador for substituído por algoritmo mais sofisticado (Voronoi, BSP), a
validação \emph{flood-fill} se torna necessária.

\subsubsection*{Telegrafia de ataques}

\textbf{GDD}: ``Cada inimigo telegrafa seus ataques com um \emph{flash} visual
antes de efetivamente agir.''\\
\textbf{Implementação}: \texttt{Bomber} pulsa visualmente durante o \emph{fuse}
($\sim$0,8\,s); \texttt{Shooter} e \texttt{Seeker} não telegrafam.\\
\textbf{Razão}: \texttt{Seeker} é \emph{melee} de contato, e a aproximação é a
própria telegrafia. \texttt{Shooter} dispara projéteis visíveis em arco amplo,
dando ao jogador 0,5 a 1,0\,s de janela de reação. Telegrafia adicional em
ambos seria redundante e poluiria a HUD.

\subsection{Adições não previstas no GDD}

\begin{itemize}
  \item \textbf{Modo de dificuldade alta}: levou a aumentar a contagem de
    inimigos, a velocidade e a agressividade do \emph{boss}.
  \item \textbf{Bússola direcional na HUD}: pequeno triângulo apontando para o
    alvo mais próximo (inimigo, ou escada quando todos estão mortos),
    adicionado para mitigar a sensação de desorientação em mapas maiores.
  \item \textbf{\emph{Mute toggle} (tecla \texttt{M})}: não previsto no GDD.
    Adicionado por boa prática de UX, evita interromper Discord, \emph{streaming}
    ou apresentação acadêmica.
  \item \textbf{Indicador de sala atual na HUD} (\texttt{SALA X/Y}): adicionado
    durante a implementação do \emph{lock-and-clear} para dar \emph{feedback} de
    progresso dentro do andar.
  \item \textbf{Tela de pausa explícita}: o GDD lista \texttt{ESC} como pausa nos
    controles, mas não detalha a UI. Implementamos um \emph{overlay}
    semi-transparente com o texto ``PAUSADO''.
\end{itemize}

\section{Cronograma: planejado vs executado}

\begin{tabular}{lll}
  \toprule
  \textbf{Marco} & \textbf{Planejado (GDD)} & \textbf{Executado} \\
  \midrule
  \emph{Setup} inicial & Semana 1 & 03 a 12 mar \\
  Sintetizador e RNG & Semana 1 & 12 mar \\
  \emph{Renderer} WebGL & Semana 2 & 22 mar \\
  Geração de \emph{dungeon} & Semana 2 & 01 abr \\
  Sistema de entidades & Semana 3 & 20 abr \\
  HUD e telas & Semana 4 & 22 a 30 abr \\
  Polimento (\emph{lock-and-clear}, \emph{spatial grid}) & não planejado & 02 a 05 maio \\
  \emph{Build} e empacotamento & Semana 4 & 07 a 09 maio \\
  Documentação & não planejado & 09 maio \\
  \bottomrule
\end{tabular}

O atraso na geração da \emph{dungeon} (esperada na semana 2, executada no fim de
março) refletiu a curva de aprendizado da geração procedural. Compensamos
comprimindo o sistema de entidades (semana 3 oficial, executado em duas semanas
e meia).

A fase de polimento (\emph{spatial grid}, \emph{lock-and-clear}, \emph{melee},
poção, pausa, bússola, ajuste de HUD) não estava no GDD original. Emergiu de
\emph{playtesting} interno e da decisão de subir a dificuldade para o modo de
dificuldade alta. Foi a fase com maior densidade de \emph{commits}: cinco em
quatro dias.

\section{Uso de IA como ferramenta de apoio}

Em pontos específicos do desenvolvimento, a equipe utilizou modelos de linguagem
como ferramenta auxiliar, da mesma forma que se usa um \emph{linter}, uma
ferramenta de \emph{static analysis} ou um colega para \emph{rubber duck
debugging}. As decisões de arquitetura, design, balanceamento e direção criativa
foram tomadas pela equipe.

Os pontos onde a IA foi consultada são bem delimitados.

\subsection{Abstrações matemáticas pontuais}

\begin{itemize}
  \item \textbf{Diferença angular normalizada} para o sistema de \emph{melee}: a
    fórmula $\mathit{diff} = \mathrm{atan2}(\sin(\alpha - \beta),
    \cos(\alpha - \beta))$ é o truque clássico para encontrar a menor diferença
    angular sem casos especiais em $\pm\pi$. A IA confirmou a derivação e ajudou
    a verificar o domínio do resultado.
  \item \textbf{Conversão HSL$\rightleftharpoons$RGB} para o \emph{tinting} do
    atlas por andar: a equação envolve interpolação não trivial nos seis
    sextantes do cone HSL. A IA forneceu a versão canônica para verificação
    cruzada.
  \item \textbf{Dimensionamento do \emph{spatial hash grid}}: a justificativa
    para células de 16 unidades, baseada na relação entre o tamanho da AABB de
    9 unidades e o custo de \emph{hashing}, foi validada como argumento, não
    copiada.
\end{itemize}

\subsection{Pesquisa de bibliotecas e \emph{trade-offs}}

\begin{itemize}
  \item Comparativo entre Electron, Tauri, Neutralino e C++ nativo, considerando
    tempos de \emph{build}, tamanho final do executável e custo de migração. A IA
    acelerou a coleta de dados que precisariam ser pesquisados em múltiplas
    fontes; a decisão de manter Electron foi da equipe, baseada no prazo e no
    requisito do GDD.
  \item Padrões consolidados de \emph{melee combat} em \emph{shooters} top-down
    (arco de \emph{hitbox} temporário com \emph{timer} curto) e de inventário
    de \emph{slot} único em \emph{roguelikes}, usados como referência durante o
    desenho da implementação.
  \item Estrutura de instrumentos do Sonant-X (parâmetros de envelope, filtro,
    LFO), consultada para entender qual combinação produz \emph{drone} ambiente
    em vez de \emph{drum} percussivo.
\end{itemize}

\subsection{Revisão de código}

A IA foi utilizada como segundo par de olhos para identificar \emph{bugs} sutis
que \emph{linters} não capturam: \emph{race conditions} entre
\texttt{setTimeout}s (\emph{gameover} contra \emph{victory}), \emph{state
leaks} entre transições de tela (teclas \emph{stuck}) e inconsistências de
\emph{textAlign} no \emph{canvas} 2D que produziam elementos da HUD fora da
tela. Cada \emph{bug} apontado foi investigado, reproduzido e corrigido
manualmente.

\subsection{Polimento da documentação}

Esta documentação \LaTeX{} foi redigida pela equipe, com auxílio da IA para
\emph{boilerplate} de configuração de pacotes, ajuste fino de \emph{layout}
(helpers \texttt{\textbackslash diagfig} e \texttt{\textbackslash
diagfiglandscape}) e revisão de estilo. O conteúdo técnico, com decisões de
arquitetura, descrição de algoritmos, \emph{trade-offs} e divergências em
relação ao GDD, reflete diretamente as escolhas documentadas no código e nas
conversas internas da equipe.

\subsection{O que a IA não fez}

Para que não haja ambiguidade: a IA não definiu o conceito do jogo, não desenhou
o sistema de combate, não escolheu a estética visual, não escreveu o GDD
original e não decidiu o ritmo dos andares. Esses são produtos do trabalho da
equipe ao longo das treze semanas de desenvolvimento.

A IA foi usada onde foi mais útil: como acelerador de pesquisa, validador de
matemática, revisor extra de código e assistente de redação técnica. Nunca como
substituto do processo de design.

\chapter{Considerações Finais}

Void Descent começou com uma proposta ambiciosa: construir um \emph{roguelike} top-down
procedural completo em JavaScript puro, sem \emph{frameworks}, sem \emph{assets}
externos e empacotado como executável \emph{desktop}. Treze semanas depois, o resultado
é um \texttt{.exe} de 73\,MB que carrega 77\,KB de código, gera mapas, áudio e
comportamentos novos a cada \emph{run}, e cumpre todos os critérios de aceite definidos
no documento de requisitos.

\section{O que foi entregue}

O escopo executado contempla:

\begin{itemize}
  \item Os oito módulos JavaScript da arquitetura, com responsabilidades isoladas e
    sem dependências circulares.
  \item Treze classes de entidade derivadas de \texttt{entity\_t}, cobrindo jogador,
    quatro tipos de inimigos, \emph{boss} de duas fases, projéteis, partículas e
    \emph{pickups}.
  \item Nove estados de jogo (\texttt{menu}, \texttt{credits}, \texttt{quit},
    \texttt{loading}, \texttt{playing}, \texttt{paused}, \texttt{shop},
    \texttt{gameover} e \texttt{victory}), todos com transições idempotentes e
    \emph{handlers} de entrada apropriados.
  \item Três andares com paletas distintas (ciano, amarelo, magenta), curva de
    dificuldade crescente e \emph{boss} no Andar 3.
  \item \emph{Lock-and-clear} de salas, \emph{spatial grid} para colisão, sistema de
    \emph{shop} com cinco \emph{upgrades} persistentes, oito armas (quatro
    \emph{ranged}, quatro \emph{melee}) e \emph{slot} de consumível.
  \item HUD funcional com vida em corações, andar, sala atual, score e combo, arma
    equipada e seu tipo, \emph{dash}, poção, contagem de inimigos, bússola direcional
    e indicador de \emph{mute}.
  \item Empacotamento Electron \emph{portable} para Windows x64, com cross-compile do
    macOS, distribuído via GitHub.
\end{itemize}

\section{O que aprendemos}

A lição mais valiosa do projeto foi que a maior parte do tempo de desenvolvimento não
está nas funcionalidades de fachada (telas, menus, HUD), mas nos detalhes invisíveis
que separam um \emph{prototype} de um produto. \emph{Race conditions} de
\texttt{setTimeout}, \emph{stuck keys} entre transições, \emph{textAlign} esquecido,
inimigos escapando de salas trancadas. Cada um desses \emph{bugs} parece pequeno
isoladamente, mas o efeito acumulado é a diferença entre algo que ``funciona'' e algo
que se sente bem de jogar.

A segunda lição foi sobre \emph{trade-offs} de empacotamento. Empacotar uma página
HTML como aplicação \emph{desktop} via Electron custa 73\,MB. Tauri ou Neutralino
custariam menos, mas com risco de incompatibilidade em máquinas antigas. C++ nativo
custaria menos ainda, mas exigiria reescrever tudo. Não há resposta universal: a
escolha certa depende do prazo, do público alvo e dos requisitos do GDD. No nosso
caso, manter Electron foi a decisão correta.

\section{O que faríamos diferente}

Em retrospecto, três pontos teriam reduzido o esforço total:

\begin{enumerate}
  \item \textbf{Spatial grid desde o início}, não como refatoração tardia. O
    \emph{loop} aninhado $O(n^2)$ funcionou enquanto havia poucos inimigos, mas o
    refator tardio tomou um dia inteiro de trabalho que poderia ter sido evitado.

  \item \textbf{Pause e \emph{lock-and-clear} planejados como parte da arquitetura}.
    Ambos foram adicionados depois e exigiram retoques em múltiplos arquivos. Tê-los
    desde o desenho inicial dos estados teria simplificado o código.

  \item \textbf{Documentação progressiva}, em vez de concentrada no fim. Esta
    documentação foi escrita em uma única semana ao final do projeto, e isso exigiu
    reconstruir o raciocínio retrospectivamente para várias decisões. Uma página de
    \emph{decision log} mantida ao longo do desenvolvimento teria capturado o porquê
    de cada escolha no momento em que foi feita.
\end{enumerate}

\section{O que vem depois}

Há um conjunto claro de melhorias que ficariam para um eventual P3 ou para uma versão
pós-disciplina:

\begin{itemize}
  \item Modos de dificuldade selecionáveis no menu. O modo de dificuldade alta
    está \emph{hardcoded} hoje; expor uma \emph{flag} permitiria oferecer um modo
    ``normal'' alinhado com o GDD original.
  \item Telegrafia visual completa para \texttt{Seekers} e \texttt{Shooters}. Hoje
    apenas \texttt{Bombers} pulsam antes de explodir.
  \item Validação \emph{flood-fill} do gerador de \emph{dungeons}. Necessária se o
    algoritmo de placement for substituído por algo mais sofisticado.
  \item Migração para Neutralino, reduzindo o \texttt{.exe} para $\sim$2\,MB. Exige
    apenas máquinas com WebView2, instalado por padrão em Windows 10 e 11
    atualizados.
  \item Tabela de \emph{high scores} local, mantida em \texttt{localStorage}.
    Reforçaria o \emph{loop} macro de competição contra o próprio histórico, sem
    quebrar o \emph{permadeath}.
\end{itemize}

\section{Encerramento}

O Void Descent existe porque quatro pessoas concordaram em entregar um jogo procedural
completo em treze semanas, e seguraram esse compromisso até o fim. O documento de
requisitos foi atendido; o GDD interno foi atendido com divergências documentadas; o
código está versionado em repositório público; o executável compila e roda.

Mais importante, ao final do processo, o jogo é divertido de jogar. Esse era o
critério verdadeiro, o que nenhum documento técnico consegue medir mas que define se
o projeto valeu a pena. Nossa avaliação interna é que valeu.


\end{document}
